\subsection{Chemistry-resolved structure of the physics tax} \label{sec:map}

Across the full test set the two arms are separated by little more than label noise. On the 3-seed scaffold split the DirectGNN control reaches an $\lnx$ MAE of $1.70 \pm 0.03$, while the grounded physics arm (learned-$\sigma$ COSMO-SAC closure) reaches $1.85 \pm 0.05$, for a global $\Delta\mathrm{MAE} = \mathrm{MAE(physics)} - \mathrm{MAE(control)} = +0.14$ (positive meaning the closure adds cost; the per-seed values span $+0.05$ to $+0.25$, so this point estimate is itself only a 3-seed mean, not a tight one; it is computed on unrounded per-seed means, so the rounded $1.70/1.85$ above differ by $0.15$). At this granularity the physics bottleneck is a small, near-noise tax on accuracy. The interesting structure appears only once the error is resolved along chemical axes, where the average tax turns out to be an average over regimes that disagree in sign.

The solvent axis is the sharp one. Binning test pairs by solvent class and recomputing $\Delta\mathrm{MAE}$ per seed (Table~\ref{tab:solvent-map}) shows the cost concentrated in solvents that hydrogen-bond strongly and directionally, precisely the regime in which the COSMO-SAC closure is least faithful and $|\lng|$ is largest. Water carries the heaviest penalty, and the effect decreases monotonically as the donor strength of the class softens. One class inverts the sign: amide solvents, which are aprotic dipolar hydrogen-bond acceptors and therefore closer to the closure's well-specified regime, are the stratum where the tax is smallest and nominally reverses sign ($\Delta\mathrm{MAE}=-0.22 \pm 0.11$ over seeds). We flag this as exploratory rather than a confirmed win: across a family of nine simultaneous solvent-class comparisons the amide interval is an uncorrected per-class pair bootstrap, and the three-seed spread alone does not exclude zero; the halogenated class shows a nominal improvement of similar magnitude but is not seed-robust. The robust reading is the \emph{direction} --- the physics tax shrinks and can turn negative as solvent hydrogen-bond donation weakens toward the acceptor-dominated regime the closure was built for --- not the significance of any single class. This solvent-resolved pattern is the chemistry-resolved shadow of the same synthetic dial we used to vary closure fidelity in controlled experiments.

\begin{table}[t]
\centering
\caption{Per-solvent-class physics tax
$\Delta\mathrm{MAE}=\mathrm{MAE(physics)}-\mathrm{MAE(control)}$ on the scaffold
split (mean $\pm$ sd over 3 seeds; positive means the closure adds cost). The
intervals are uncorrected per-class pair bootstraps across a family of nine
simultaneous comparisons; the two negative (physics-helps) entries are
exploratory, not multiplicity-controlled.}
\label{tab:solvent-map}
\begin{tabular}{lr}
\toprule
Solvent class & $\Delta\mathrm{MAE}$ (mean $\pm$ sd) \\
\midrule
Water & $+0.52 \pm 0.35$ \\
Carboxylic acid & $+0.39 \pm 0.16$ \\
Sulfoxide & $+0.26 \pm 0.24$ \\
Nitrile & $+0.21 \pm 0.19$ \\
Alcohol & $+0.15 \pm 0.14$ \\
Aromatic & $+0.03 \pm 0.63$ \\
Hydrocarbon & $+0.02 \pm 0.50$ \\
Halogenated & $-0.29 \pm 0.36$ \\
Amide & $-0.22 \pm 0.11$ \\
\bottomrule
\end{tabular}
\end{table}

The solute axis tells a compatible but weaker story in prose. Oxygenated solutes ($+0.32$), polyaromatics ($+0.27$), and heterocycles ($+0.17$) are where the closure costs the most, while halogenated-aromatic solutes are roughly neutral. An earlier hypothesis, that the physics path should hurt most on charged and high-$|\lng|$ solutes such as salts, did not replicate once the split was corrected: the charged/salt stratum is neutral at $-0.04 \pm 0.19$. The clean penalty falls instead on ordinary neutral organics, not on the ionic tail we had expected to be hardest.

The most suggestive pattern is a novelty inversion that runs opposite to the usual coverage intuition. Binning solutes into three pre-specified nearest-train Tanimoto-similarity strata, the physics arm is neutral on the most novel chemistry ($\leq 0.3$: $-0.04 \pm 0.18$) and hurts most on the most familiar chemistry ($>0.8$: $+0.51 \pm 0.09$). We flag this as exploratory, not a confirmed effect: the $\pm0.09$ is the spread across three training seeds on the fixed test split, not a within-bin sampling interval over which solutes fall in the stratum, and the Tanimoto binning is one of several strata families we examine (solvent class, solute class, similarity) without multiplicity control. Read as a direction rather than a tested quantity, it suggests the ceiling on the physics path is not coverage or extrapolation but the closure's systematic misfit---masked on novel solutes by the large errors both arms already incur there, and exposed on easy, well-covered solutes where the direct control is accurate enough for a structural bias in the closure to become the dominant residual---consistent with the same closure-misspecification reading as the rest of the paper. A multiplicity-controlled, within-bin-bootstrapped test is left to future work. Read together, the maps say the physics decomposition pays where the closure is misspecified relative to the solvent chemistry, and it can only turn into a win in the narrow acceptor-dominated regime the closure was built to describe.
